\pagenumbering{roman}
\setcounter{page}{1}

%=======================================================================
\thispagestyle{empty}
\begin{center}
\singlespacing
\vspace*{0.5in}

{\large\bfseries
DISCRETE-TIME NEURAL NETWORK BASED STATE OBSERVER\\[0.4em]
WITH COMMAND-FILTERED NEURAL NETWORK CONTROL FOR A CLASS\\[0.4em]
OF SYSTEMS WITH UNMATCHED UNCERTAINTIES}

\vspace{0.8in}
by
\vspace{0.4in}

JASON MICHAEL STUMFOLL

\vspace{0.8in}
A THESIS
\vspace{0.4in}

Presented to the Faculty of the Graduate School of the

\vspace{0.3in}
MISSOURI UNIVERSITY OF SCIENCE AND TECHNOLOGY

\vspace{0.4in}
In Partial Fulfillment of the Requirements for the Degree

\vspace{0.5in}
MASTER OF SCIENCE IN AEROSPACE ENGINEERING

\vspace{0.5in}
Revision 8

\vfill
\end{center}
\doublespacing
\clearpage

%=======================================================================
\chapter*{ABSTRACT}
\addcontentsline{toc}{chapter}{ABSTRACT}

An observer is a dynamic system that estimates the state variables of another
system from noisy measurements, either to reconstruct unmeasured states or to
improve the accuracy of measured ones. The Modified State Observer (MSO) is a
standard observer structure augmented with an online neural network that
estimates both the system states and the system uncertainty. It has been applied
to orbit uncertainty estimation, atmospheric reentry uncertainty estimation, and
the control of a nonlinear electrohydraulic system with parameter uncertainty.

This work develops the discrete-time extension of the MSO (DMSO) and establishes
its stability using Lyapunov theory. \rev{The stability argument presented here
evaluates the Lyapunov first difference exactly rather than bounding it term by
term. Doing so exposes a family of cancellations produced by the adaptation law
itself, and yields admissibility conditions that are both satisfiable and
physically interpretable: the error-injection matrix must satisfy
$\sigmax(F-\Km)<1$, and the adaptation gain must be bounded above. In
particular, small adaptation gains are admissible, which a term-by-term bound
incorrectly excludes.}

To demonstrate the DMSO, simulation studies are performed on a two-wheeled
inverted pendulum (TWIP) robot, an unstable nonlinear system with unmatched
uncertainties, and the observer is implemented on hardware.

\rev{A command-filtered backstepping controller with neural network uncertainty
compensation is then developed for the TWIP. The controller uses the tilt angle
as a virtual control for the position subsystem. Relative to a two-step design,
the formulation here (i) retains the coupling term between the position and
attitude error subsystems and cancels it explicitly, (ii) assigns damping to
every error coordinate rather than only to the two driven coordinates, and (iii)
uses command filtering so that the functions approximated by the neural networks
depend only on measured states and the reference trajectory, removing the
algebraic dependence of each network's target on the other network's output.
Uniform ultimate boundedness of all tracking errors and all weight estimation
errors follows, with $\sigma$-modification supplying weight boundedness without
a persistency-of-excitation requirement.}

\clearpage

%=======================================================================
\chapter*{ACKNOWLEDGMENTS}
\addcontentsline{toc}{chapter}{ACKNOWLEDGMENTS}

I want to acknowledge my advisor, Dr.\ S.\ N.\ Balakrishnan, for all of his
support and guidance throughout this project. I want to thank my committee,
Dr.\ Jagannathan Sarangapani and Dr.\ Joshua Rovey, for their classes and for
the research opportunities provided throughout my collegiate career. A large
part of this goes to Avimanyu Sahoo, who proved invaluable during my
difficulties with the mathematics necessary for this work. I also want to thank
my parents for their support through my years in school, and, lastly but not
least, the love of my life, for being there with me every step of the way.

\clearpage

%=======================================================================
\singlespacing
\tableofcontents
\clearpage
\listoffigures
\addcontentsline{toc}{chapter}{LIST OF ILLUSTRATIONS}
\clearpage
\listoftables
\addcontentsline{toc}{chapter}{LIST OF TABLES}
\clearpage

%=======================================================================
\chapter*{NOMENCLATURE}
\addcontentsline{toc}{chapter}{NOMENCLATURE}

\begin{tabbing}
\hspace{1.3in}\=\kill
$x_k$        \> System state vector at time step $k$, $x_k\in\R^n$ \\
$\xhat_k$    \> State estimate \\
$e_k$        \> State estimation error, $e_k=x_k-\xhat_k$ \\
$F,\,G$      \> Discretized system and input matrices \\
$\Bm$        \> Uncertainty selection matrix, $\Bm\in\R^{n\times p}$ \\
$\Km$        \> Observer error-injection gain \\
$A$          \> Error-injection matrix, $A=F-\Km$ \\
$f(\cdot)$   \> Unknown uncertainty vector \\
$W,\,\What_k$\> Ideal and estimated neural network weights \\
$\Wt_k$      \> Weight estimation error, $\Wt_k=W-\What_k$ \\
$\phi_k$     \> Basis (activation) vector, $\phi_k=\phi(x_k)$ \\
$\psi_k$     \> Functional approximation error, $\psi_k=\Wt_k^T\phi_k$ \\
$\epsilon_k$ \> Network reconstruction error, $\norm{\epsilon_k}\le\epsilon_N$ \\
$\Gamma$     \> Observer network adaptation gain \\
$\kappa$     \> $\sigma$-modification (leakage) gain \\
$z_i$        \> Backstepping error coordinates, $i=1,\dots,4$ \\
$\chi_i$     \> Command filter errors \\
$c_i$        \> Backstepping damping gains \\
$\varpi_i$   \> Command filter bandwidths \\
$\theta$     \> Pendulum tilt angle from vertical \\
$V$          \> Motor terminal voltage (control input) \\
\end{tabbing}
\doublespacing
\clearpage
